\section{程序实现与试验设计}
\label{sec:experiments}

程序采用分层结构：
\texttt{core} 负责电网对象、导纳矩阵和功率方程；
\texttt{solvers} 提供统一求解接口；
\texttt{cpf} 负责专门对连续潮流进行预测、校正、参数化和步长控制；
\texttt{analysis} 比较分析算法、鲁棒性、临界停滞与规模增长。

为保证结果可比较，所有算法使用相同的输入数据、初值规则和残差定义；
重复计时先执行一次预热，再按单侧 Tukey--IQR 规则识别异常耗时。
临界负荷测试开启发电机无功限值及 PV--PQ 转换，
并分别记录“方程残差满足阈值”和“算法停止但残差非零”两类状态。

\begin{table}[H]
\centering
\caption{主要试验设置}
\label{tab:settings}
\begin{tabularx}{\textwidth}{p{3.8cm}X}
\toprule
项目 & 设置\\
\midrule
测试系统 & IEEE 3、4、5、9、14、30、39、43、57、118、145、162 和 300 节点\\
基础条件 & 100 MVA；平坦启动；基础收敛精度 $10^{-6}$\\
重复计时 & 每种算法运行 11 次，第 1 次预热，正式样本采用单侧 IQR 规则识别异常值\\
精度扫描 & $10^{-4}$、$10^{-6}$、$10^{-8}$、$10^{-10}$\\
初值扰动 & 电压幅值 $\pm2\%$、$\pm5\%$、$\pm10\%$；相角 $\pm2^\circ$、$\pm5^\circ$、$\pm10^\circ$\\
病态工况 & 负荷倍率 1.600--1.620；线路电阻倍率 1、2、3；开启 PV--PQ 转换\\
\bottomrule
\end{tabularx}
\end{table}
